% §8.4 The Bloom Filter's Gap
%
% Monograph register: the log2(e) ~ 1.44 overhead of the standard Bloom filter
% over the entropy floor, the gap table, and the structural reason. rem:compressed-bloom
% salvages ex:ch8-10 with the CORRECTED Mitzenmacher mechanism: the half-full
% optimal array is incompressible (max per-bit entropy); compression helps only a
% SPARSER array, approaching the floor as eps->0 without reaching it, and the
% residual gap is structural. Cross-refs def:bloom-filter, thm:bloom-optimal-k,
% eq:bloom-optimal-fpr, eq:bloom-fpr, thm:entropy-floor, eq:entropy-floor,
% ch:optimal-bernoulli-sets, sec:ch8-what-comes-next. DEFINED: sec:bloom-gap,
% prop:bloom-overhead, eq:bloom-optimal-fpr-recall, eq:bloom-bits-per-element,
% rem:compressed-bloom.

\section{The Bloom Filter's Gap}
\label{sec:bloom-gap}

The standard Bloom filter is FNR-zero (\Cref{def:bloom-filter}), so
\Cref{thm:entropy-floor} governs it. How far above the floor does it sit? The
answer is a single constant, $\log_2 e \approx 1.443$, the same at every
$\fprate$.

\subsection*{Bits per element}

By \Cref{thm:bloom-optimal-k} the optimal Bloom filter is half-occupied with
\begin{equation}
\label{eq:bloom-optimal-fpr-recall}
  \fprate^* \;=\; 2^{-(m/n)\ln 2}.
\end{equation}
Inverting at $\fprate^* = \fprate$, $\log_2 \fprate = -(m/n)\ln 2$, so
\begin{equation}
\label{eq:bloom-bits-per-element}
  \frac{m}{n} \;=\; \frac{\log_2(1/\fprate)}{\ln 2}
  \;=\; \log_2(1/\fprate)\cdot \log_2 e,
\end{equation}
using $1/\ln 2 = \log_2 e$. The Bloom budget is thus a constant multiple
$\log_2 e$ of the floor cost $\log_2(1/\fprate)$ of \Cref{thm:entropy-floor},
so the ratio is $\log_2 e = 1/\ln 2 \approx 1.4427$, independent of $\fprate$.

\begin{proposition}[Bloom filter overhead]
\label{prop:bloom-overhead}
At every $\fprate \in (0,1)$ the optimal standard Bloom filter uses $m/n =
\log_2(1/\fprate)\cdot\log_2 e$ bits per element, which is $\log_2 e \approx
1.4427$ times the floor of \Cref{thm:entropy-floor}. The relative overhead is
constant.
\end{proposition}

The constant is the chapter-5 half-occupancy result read dually: $\ln 2 =
1/\log_2 e$ is the optimal load that maximizes per-bit entropy, and the same
$\log_2 e$ is the gap to the floor. Half-occupancy is information-theoretically
the wrong target, but it is the best a representation built from $k$
independent hashes and a flat bit array can do.

\subsection*{The gap in absolute terms}

A constant factor on a function that grows like $\log_2(1/\fprate)$ produces an
absolute gap that grows with it:
\begin{center}
\begin{tabular}{l|rrr}
\hline
  $\fprate$ & floor (bits/el) & Bloom (bits/el) & gap (bits/el) \\
\hline
  $10^{-2}$ & $6.64$  & $9.59$  & $2.95$ \\
  $10^{-3}$ & $9.97$  & $14.38$ & $4.41$ \\
  $10^{-6}$ & $19.93$ & $28.76$ & $8.83$ \\
  $10^{-9}$ & $29.90$ & $43.14$ & $13.25$ \\
\hline
\end{tabular}
\end{center}
At one-percent FPR the filter pays about three extra bits per element; at
one-in-a-million, about nine. For a million-element set at $\fprate = 10^{-6}$
the gap is $8.83$ megabits, about $1.10$ megabytes, that the standard
construction cannot recover. A constant factor on an already-small quantity is
acceptable in many settings, and Bloom filters dominate the ecosystem for good
reason; the point of \Cref{prop:bloom-overhead} is to quantify the slack, not
to condemn the construction, and to ask which design choice forces it.

\subsection*{Why the gap is structural}

The Bloom filter pays its overhead because of how it is built: $k$ independent
hashes write into a single flat bit array, and a query accepts iff all $k$ read
bits are $1$. The optimal-$k$ analysis sets the per-bit fill to $p = 1/2$,
where the array carries one full bit of Shannon entropy per array bit and the
encoding is maximally dense at the bit level. But the floor is set by the
membership problem, not by the encoding's internal density. The Bloom filter
maximizes per-bit entropy and takes the $k$th power of it, getting a clean
optimum at half-occupancy while spending bits on the array's internal balance
rather than on the discrimination it was meant to perform.

\begin{remark}[Compressed Bloom filters]
\label{rem:compressed-bloom}
Half-occupancy is exactly what makes the optimal array \emph{incompressible}:
at $p = 1/2$ each bit carries a full bit of entropy, so arithmetic coding of
the optimal-$k$ array recovers nothing. \textcite{mitzenmacher-2002} therefore
goes the other way, using a \emph{sparser} array (fewer hashes, fill $p \neq
1/2$) whose low per-bit entropy compresses well, trading a larger in-memory
array and the loss of in-place random-access queries for a smaller transmitted
size. The compressed cost per element is $\log_2(1/\fprate)\cdot g(p)$ with
$g(p) = \frac{-p\ln p - (1-p)\ln(1-p)}{\ln p\,\ln(1-p)}$: at $p = 1/2$,
$g = 1/\ln 2 \approx 1.44$ (the uncompressed overhead, no gain), while $g(p)
\to 1$ as $p \to 0$, so the sparse limit \emph{approaches} the floor as
$\fprate \to 0$ without reaching it at finite $\fprate$. Compression recovers
the per-bit redundancy of a non-half-full array; the residual gap is
structural, and closing it requires abandoning the flat-array-with-$k$-hashes
representation entirely.
\end{remark}

That is what chapter~14 does. The Bernoulli Hash Function of
\Cref{ch:optimal-bernoulli-sets} replaces the flat array with a
perfect-hash-based encoding: each element gets a deterministic location from a
perfect hash, and a fingerprint of $\log_2(1/\fprate)$ bits is stored there, so
collisions in the fingerprint table are the only source of false positives and
the $\log_2 e$ overhead vanishes asymptotically. The gap is not a deficiency of
analysis but a consequence of a specific structural choice, and a different
choice closes it. \Cref{sec:ch8-what-comes-next} sketches the route there.
